\documentclass[12pt,a4paper]{report}

\usepackage[a4paper,margin=1in]{geometry}
\usepackage[T1]{fontenc}
\usepackage{setspace}
\usepackage{graphicx}
\usepackage{array}
\usepackage{booktabs}
\usepackage{longtable}
\usepackage{float}
\usepackage{hyperref}
\usepackage{enumitem}
\usepackage{caption}
\usepackage{fancyhdr}
\usepackage{titlesec}

\setstretch{1.2}
\setlength{\parskip}{6pt}
\setlength{\parindent}{0pt}
\setlength{\headheight}{14.5pt}

\pagestyle{fancy}
\fancyhf{}
\fancyhead[L]{\nouppercase{\leftmark}}
\fancyhead[R]{Mini Project Report}
\fancyfoot[C]{\thepage}

\titleformat{\chapter}[display]
{\bfseries\Huge}
{Chapter \thechapter}
{0.5ex}
{\Large}

\graphicspath{{figures/}{./}}
\captionsetup[figure]{justification=centering,singlelinecheck=false}

\newcommand{\screenshotslot}[2]{
\begin{figure}[!htbp]
\centering
\includegraphics[width=0.92\textwidth,height=0.62\textheight,keepaspectratio]{#1}
\caption{#2}
\end{figure}
}

\begin{document}

% ------------------------- TITLE PAGE -------------------------
\begin{titlepage}
\centering

\vspace*{1cm}
{\Large \textbf{A MINI PROJECT REPORT ON}}\\[0.6cm]
{\Huge \textbf{RetailFlow Business Portal}}\\[0.6cm]
{\Large \textbf{(Dynamic Web Application for Business Functionality)}}\\[1.2cm]

{\large \textbf{SUBMITTED TO}}\\[0.3cm]
{\Large \textbf{SAVITRIBAI PHULE PUNE UNIVERSITY, PUNE}}\\[1.2cm]

{\large \textbf{SUBMITTED BY}}\\[0.4cm]
\begin{tabular}{p{8cm}p{5cm}}
\textbf{Vishwajeet Tulse} & \hspace{1.7cm}\textbf{Roll No: T5543}  \\
\end{tabular}

\vfill
\includegraphics[width=6cm,height=4cm,keepaspectratio]{\detokenize{SinhgadLogo (1).jpg}}
\vfill

{\large \textbf{UNDER THE GUIDANCE OF}}\\[0.3cm]
{\large Dr. S. P. Bendale}\\[1cm]

{\large \textbf{DEPARTMENT OF COMPUTER ENGINEERING}}\\
{\large \textbf{NBN SINHGAD TECHNICAL INSTITUTES CAMPUS}}\\
{\large \textbf{NBN SINHGAD SCHOOL OF ENGINEERING}}\\
{\large 10/1, Ambegaon (Bk), Pune - 411041}\\[0.8cm]

{\large \textbf{Academic Year 2025--2026}}

\end{titlepage}

\pagenumbering{roman}

% ------------------------- CERTIFICATE -------------------------
\chapter*{}
\addcontentsline{toc}{chapter}{CERTIFICATE}
\vspace*{-1.5cm}
\begin{center}
{\Huge\bfseries CERTIFICATE\par}
\vspace{1cm}
\includegraphics[width=6cm,height=4cm,keepaspectratio]{\detokenize{SinhgadLogo (1).jpg}}
\end{center}

This is to certify that the project report entitled \textbf{RetailFlow Business Portal}
submitted by

\noindent
\begin{tabular*}{\textwidth}{@{\extracolsep{\fill}}l r}
\textbf{Vishwajeet Tulse} & \textbf{Roll No: T5543} \\
\end{tabular*}

is a bonafide work carried out under our supervision and guidance.
This work is submitted in partial fulfillment of the requirements of
Web Technology Laboratory of Savitribai Phule Pune University, Pune.

\vspace{3cm}
\noindent
\begin{minipage}[t]{0.48\textwidth}
\raggedright
Guide Signature\\
\textbf{Dr. S. P. Bendale}
\end{minipage}
\hfill
\begin{minipage}[t]{0.48\textwidth}
\raggedleft
HOD Signature\\
\textbf{Dr. S. P. Bendale}
\end{minipage}

\vspace{1.5cm}
\textbf{Date:} \rule{3cm}{0.15mm} \hfill
\textbf{Place:} Pune

% ------------------------- ACKNOWLEDGEMENT -------------------------
\chapter*{ACKNOWLEDGEMENT}
\addcontentsline{toc}{chapter}{ACKNOWLEDGEMENT}

It gives me great pleasure to present this project report on
\textbf{Movie Recommendation Model using Scikit-learn}. I express sincere gratitude
and thanks to my guide \textbf{Dr. S. P. Bendale} for constant guidance, motivation,
and valuable suggestions throughout the mini project work.

I am thankful to the Head of Department, all faculty members, and the institute
for providing the required infrastructure and support. I also thank my family and
friends for their encouragement and support during the completion of this project.

% ------------------------- ABSTRACT -------------------------
\chapter*{ABSTRACT}
\addcontentsline{toc}{chapter}{ABSTRACT}

This mini project presents a \textbf{movie recommendation model} implemented in
Python using the \textbf{scikit-learn} library. The system uses the
\texttt{movie\_dataset.csv} dataset from the GitHub repository
\href{https://github.com/rashida048/Some-NLP-Projects}{Some-NLP-Projects}.

The project follows a \textbf{content-based recommendation} approach. Important text
features such as genres, keywords, cast, director, tagline, and overview are
combined and transformed into numerical vectors using
\textbf{TF-IDF (Term Frequency - Inverse Document Frequency)}. Similarity between
movies is then computed using \textbf{cosine similarity}.

Given an input movie title, the system returns top-N similar movies. The entire
workflow is implemented in a Jupyter notebook for easy demonstration,
experimentation, and analysis. The project demonstrates practical use of machine
learning concepts, text preprocessing, feature engineering, and recommendation
techniques in a real-world dataset.

% ------------------------- CONTENTS -------------------------
\tableofcontents
\clearpage

\pagenumbering{arabic}

% ------------------------- CHAPTER 1 -------------------------
\chapter{Introduction}

\section{Project Title}
\textbf{Develop a Movie Recommendation Model Using Scikit-learn}

\section{Introduction}
Recommendation systems are widely used in digital platforms to help users discover
relevant items from large collections. In the movie domain, recommending similar
films improves user experience and reduces search effort.

In this project, a movie recommendation model is developed using Python and
scikit-learn. The model uses movie metadata and applies natural language feature
engineering to identify similarity among movies. The implementation is completed in
Jupyter notebook format and validated on a publicly available dataset.

\section{Objectives of the Project}
\begin{itemize}[leftmargin=1.2cm]
\item To build a movie recommendation system using scikit-learn.
\item To preprocess and combine multiple textual metadata fields.
\item To apply TF-IDF vectorization for feature representation.
\item To compute movie-to-movie similarity using cosine similarity.
\item To return top-N recommendations for a given movie title.
\end{itemize}

\section{Scope}
The current project focuses on a content-based recommender using one dataset and
single-user query input. Future scope includes hybrid recommendation (content +
collaborative filtering), user profiling, and web deployment.

% ------------------------- CHAPTER 2 -------------------------
\chapter{Literature Review}

Recommendation systems are generally categorized into collaborative filtering,
content-based filtering, and hybrid approaches.

Collaborative filtering depends on user-item interaction data such as ratings.
Although powerful, it often suffers from cold-start issues when new users or items
have little history.

Content-based recommendation uses item attributes and recommends items similar to
those preferred by the user. For text-rich movie metadata, TF-IDF with cosine
similarity is a widely accepted baseline due to simplicity, interpretability, and
good practical performance.

Based on these observations, this mini project adopts a content-based method using
movie metadata and scikit-learn utilities.

% ------------------------- CHAPTER 3 -------------------------
\chapter{Software Requirement Specification}

\section{Functional Requirements}
\begin{itemize}[leftmargin=1.2cm]
\item Load movie metadata from \texttt{data/movie\_dataset.csv}.
\item Clean and preprocess text fields (lowercasing, symbol removal, spacing).
\item Build a TF-IDF feature matrix from selected metadata columns.
\item Accept movie title input and support exact/partial title match.
\item Compute similarity and display top-N recommended movies.
\item Provide notebook-based execution for demonstration and testing.
\end{itemize}

\section{Non-Functional Requirements}
\begin{itemize}[leftmargin=1.2cm]
\item Usability: Simple notebook flow with clear outputs.
\item Performance: Fast recommendation on medium-size dataset.
\item Maintainability: Modular functions for preprocessing and recommendation.
\item Reliability: Error handling for missing files and unknown movie titles.
\item Reproducibility: Same results for same input and dataset.
\end{itemize}

\section{Hardware Requirements}
\begin{itemize}[leftmargin=1.2cm]
\item Processor: Intel Core i3 or above
\item RAM: 4 GB or above (8 GB recommended)
\item Storage: 1 GB free disk space
\end{itemize}

\section{Software Requirements}
\begin{itemize}[leftmargin=1.2cm]
\item Operating System: Windows 10/11 or Linux
\item Python: 3.10 or above
\item Notebook Environment: Jupyter Notebook / VS Code Notebook
\item Libraries: pandas, scikit-learn
\item Optional IDE: VS Code
\end{itemize}

% ------------------------- CHAPTER 4 -------------------------
\chapter{System Analysis and Design}

\section{Existing System}
Traditional movie browsing is keyword-based and often not personalized. Users may
need to search multiple titles manually before finding relevant content.

\section{Proposed System}
The proposed model recommends movies similar to a selected title based on metadata.
Key advantages are:
\begin{itemize}[leftmargin=1.2cm]
\item No dependency on user rating history
\item Efficient similarity search using vector space representation
\item Easy implementation and explainable recommendation logic
\end{itemize}

\section{Dataset Overview}
Dataset file used:
\begin{itemize}[leftmargin=1.2cm]
\item \texttt{movie\_dataset.csv}
\end{itemize}

Important fields considered:
\begin{itemize}[leftmargin=1.2cm]
\item \texttt{title}
\item \texttt{genres}
\item \texttt{keywords}
\item \texttt{tagline}
\item \texttt{cast}
\item \texttt{director}
\item \texttt{overview}
\end{itemize}

\section{Recommendation Model}
Text from selected columns is combined into one corpus per movie.
TF-IDF vectors are generated using scikit-learn. Similarity is measured using
cosine similarity.

TF-IDF scoring formula:
\[
\mathrm{tfidf}(t,d) = \mathrm{tf}(t,d) \times \log\left(\frac{N}{\mathrm{df}(t)}\right)
\]

Cosine similarity formula:
\[
\mathrm{sim}(A,B) = \frac{A \cdot B}{\|A\|\|B\|}
\]

Movies with highest similarity scores are returned as recommendations.

% ------------------------- CHAPTER 5 -------------------------
\chapter{System Design and Implementation}

\section{Implementation Workflow}
The notebook implementation follows this sequence:
\begin{itemize}[leftmargin=1.2cm]
\item Import required Python libraries
\item Load dataset from local path
\item Preprocess and clean text features
\item Generate TF-IDF feature matrix
\item Define recommendation function
\item Run sample and interactive recommendations
\end{itemize}

\section{Notebook and Project Files}
\begin{itemize}[leftmargin=1.2cm]
\item \texttt{movie\_recommendation\_project.ipynb} - Main implementation
\item \texttt{data/movie\_dataset.csv} - Dataset
\item \texttt{requirements.txt} - Python dependencies
\end{itemize}

\section{Algorithmic Highlights}
\begin{itemize}[leftmargin=1.2cm]
\item Combined multiple textual attributes for richer context
\item Used bi-gram TF-IDF settings for better phrase representation
\item Added partial title matching for user-friendly query support
\item Returned ranked recommendations with similarity scores
\end{itemize}

\section{Sample Output}
For input movie title \textbf{Avatar}, the model returns top similar movies such as:
\textit{Aliens, Cargo, Star Trek Beyond, Guardians of the Galaxy, Alien}, etc.

% ------------------------- CHAPTER 6 -------------------------
\chapter{Testing}

\section{Testing Techniques Used}
\begin{itemize}[leftmargin=1.2cm]
\item Functional Testing: Validation of each notebook step
\item Data Validation Testing: File existence and column checks
\item Result Verification: Relevance of recommended movie list
\item Error Handling Testing: Unknown title and missing file scenarios
\end{itemize}

\section{Sample Test Cases}
\begin{table}[H]
\centering
\begin{tabular}{p{1.7cm}p{3.6cm}p{3.2cm}p{4.0cm}p{1.8cm}}
\toprule
\textbf{TC ID} & \textbf{Function} & \textbf{Input} & \textbf{Expected Output} & \textbf{Status} \\
\midrule
TC01 & Dataset Load & Valid CSV file path & DataFrame loaded successfully & Pass \\
TC02 & Preprocessing & Text columns with nulls & Clean combined corpus generated & Pass \\
TC03 & TF-IDF Build & Clean corpus & TF-IDF matrix created & Pass \\
TC04 & Exact Search & Avatar & Top-N recommendations displayed & Pass \\
TC05 & Partial Search & Batman & Closest title selected, results shown & Pass \\
TC06 & Invalid Movie & RandomXYZ123 & Error message: Movie not found & Pass \\
TC07 & Missing Dataset & Deleted CSV file & FileNotFoundError displayed & Pass \\
TC08 & Output Quality & Top-N = 10 & Ranked list with similarity values & Pass \\
\bottomrule
\end{tabular}
\end{table}

% ------------------------- CHAPTER 7 -------------------------
\chapter{Conclusion}

The mini project successfully implements a movie recommendation model using
scikit-learn and real-world movie metadata. The model applies content-based
filtering through TF-IDF feature extraction and cosine similarity scoring.

The developed notebook is simple, modular, and easy to demonstrate. It satisfies
the project objective of building a recommender in Python while also showing
practical understanding of text preprocessing, vectorization, and ranking.

This project can be extended with hybrid recommendation methods, user feedback,
web-based interfaces, and formal evaluation metrics such as Precision@K and
Recall@K.

\end{document}
